% Microeconomics*
% Final Exam covering Chs. 1-4, 6-9
%----------------------------------------------
\documentclass[12pt]{exam}
 
\usepackage{setspace}
\setstretch{1.15}
\nopointsinmargin
\pointformat{}

\usepackage{problemset}

\usepackage[margin=1in]{geometry}

%\noprintanswers
\printanswers

\begin{document}

\FontEasyRead

\Heading
{Microeconomics*}
{Cumulative Test}
{75 Minutes, 100 points plus 10 bonus points.
    Annotated "cheat-sheet" permitted; no other resources
    Answer all multiple-choice questions by selecting the best answer.
    Answer all short-response questions clearly, using correct reasoning.
    Show all work for math problems.}

%%%%%%%%%%%%%%%%%%%%%%%%%%%%%%%%%%%%%%%%%%%%%%%%%%%%%%
\section*{Part I: Multiple Choice (2 points each)}

\begin{questions}

    \begin{mcquestion}[2]
        Opportunity cost is best defined as:
        \begin{choices}
            \choice The monetary cost of a decision only
            \CorrectChoice The value of the next best alternative foregone
            \choice Total accounting profit
            \choice The price paid for a good
        \end{choices}
    \end{mcquestion}

    \begin{mcquestion}[2]
        The law of demand states that:
        \begin{choices}
            \choice As price increases, quantity demanded increases
            \CorrectChoice As price increases, quantity demanded decreases
            \choice Demand is unrelated to price
            \choice Supply determines demand
        \end{choices}
    \end{mcquestion}

    \begin{mcquestion}[2]
        A market is in equilibrium when:
        \begin{choices}
            \choice Quantity demanded exceeds quantity supplied
            \choice Price is changing rapidly
            \CorrectChoice Quantity demanded equals quantity supplied
            \choice Firms earn zero revenue
        \end{choices}
    \end{mcquestion}

    \begin{mcquestion}[2]
        If income increases and a good is normal, demand will:
        \begin{choices}
            \choice Decrease
            \choice Stay the same
            \CorrectChoice Increase
            \choice Become perfectly inelastic
        \end{choices}
    \end{mcquestion}

    \begin{mcquestion}[2]
        Price elasticity of demand measures:
        \begin{choices}
            \choice How price responds to supply changes
            \CorrectChoice Responsiveness of quantity demanded to price changes
            \choice Total revenue changes only
            \choice Cost changes in production
        \end{choices}
    \end{mcquestion}

    \begin{mcquestion}[2]
        If demand is elastic, a price increase will:
        \begin{choices}
            \choice Increase total revenue
            \CorrectChoice Decrease total revenue
            \choice Leave revenue unchanged
            \choice Eliminate demand
        \end{choices}
    \end{mcquestion}

    \begin{mcquestion}[2]
        Diminishing marginal returns occurs when:
        \begin{choices}
            \choice Output increases proportionally with inputs
            \CorrectChoice Each additional unit of input adds less output than the previous
            \choice Costs decrease as output increases
            \choice Firms shut down
        \end{choices}
    \end{mcquestion}

    \begin{mcquestion}[2]
        Fixed costs are:
        \begin{choices}
            \choice Costs that vary with output
            \CorrectChoice Costs that do not vary with output
            \choice Always zero in the long run
            \choice Equal to marginal cost
        \end{choices}
    \end{mcquestion}

    \begin{mcquestion}[2]
        Marginal cost is:
        \begin{choices}
            \choice Total cost divided by quantity
            \CorrectChoice The cost of producing one additional unit
            \choice Fixed cost per unit
            \choice Average revenue
        \end{choices}
    \end{mcquestion}

    \begin{mcquestion}[2]
        A profit-maximizing firm in perfect competition produces where:
        \begin{choices}
            \choice Price equals total cost
            \choice Marginal revenue is zero
            \CorrectChoice Price equals marginal cost
            \choice Average cost is minimized
        \end{choices}
    \end{mcquestion}

    \begin{mcquestion}[2]
        In the short run, a firm should shut down if:
        \begin{choices}
            \choice Price is below average total cost
            \CorrectChoice Price is below average variable cost
            \choice Revenue is positive
            \choice Marginal cost equals price
        \end{choices}
    \end{mcquestion}

    \begin{mcquestion}[2]
        Consumer surplus is:
        \begin{choices}
            \choice Total revenue minus total cost
            \CorrectChoice Willingness to pay minus actual price
            \choice Firm profit
            \choice Tax revenue
        \end{choices}
    \end{mcquestion}

    \begin{mcquestion}[2]
        Producer surplus is:
        \begin{choices}
            \choice Price minus marginal utility
            \CorrectChoice Price minus marginal cost (area above supply curve)
            \choice Total cost minus revenue
            \choice Consumer willingness to pay
        \end{choices}
    \end{mcquestion}

    \begin{mcquestion}[2]
        A binding price ceiling causes:
        \begin{choices}
            \choice Surplus
            \CorrectChoice Shortage
            \choice No effect
            \choice Increased equilibrium price
        \end{choices}
    \end{mcquestion}

    \begin{mcquestion}[2]
        A binding price floor causes:
        \begin{choices}
            \choice Shortage
            \CorrectChoice Surplus
            \choice Market equilibrium
            \choice Zero production
        \end{choices}
    \end{mcquestion}

    \begin{mcquestion}[2]
        If supply increases, equilibrium price will:
        \begin{choices}
            \choice Increase
            \CorrectChoice Decrease
            \choice Stay constant
            \choice Become infinite
        \end{choices}
    \end{mcquestion}

    \begin{mcquestion}[2]
        If demand is inelastic, consumers tend to:
        \begin{choices}
            \choice Avoid price changes entirely
            \CorrectChoice Bear more of the tax burden
            \choice Shift supply curves
            \choice Eliminate scarcity
        \end{choices}
    \end{mcquestion}

    \begin{mcquestion}[2]
        Long-run competitive equilibrium results in:
        \begin{choices}
            \choice Positive economic profit
            \CorrectChoice Zero economic profit
            \choice Monopoly pricing
            \choice Excess demand
        \end{choices}
    \end{mcquestion}

    \begin{mcquestion}[2]
        Economic profit differs from accounting profit because it:
        \begin{choices}
            \choice Ignores revenue
            \CorrectChoice Includes opportunity costs
            \choice Ignores costs
            \choice Equals total revenue
        \end{choices}
    \end{mcquestion}

    \begin{mcquestion}[2]
        A firm is a price taker when:
        \begin{choices}
            \choice It sets market price
            \CorrectChoice It cannot influence market price
            \choice It is a monopoly
            \choice It has no costs
        \end{choices}
    \end{mcquestion}

    \begin{mcquestion}[2]
        The shutdown point occurs when:
        \begin{choices}
            \choice Profit is zero
            \CorrectChoice Price equals minimum AVC
            \choice Revenue exceeds cost
            \choice Marginal revenue is highest
        \end{choices}
    \end{mcquestion}

    \begin{mcquestion}[2]
        Total revenue is:
        \begin{choices}
            \choice Price minus quantity
            \CorrectChoice Price times quantity
            \choice Marginal cost times quantity
            \choice Fixed cost plus variable cost
        \end{choices}
    \end{mcquestion}

    \begin{mcquestion}[2]
        A shift in supply is caused by:
        \begin{choices}
            \choice A change in price
            \CorrectChoice A change in input costs
            \choice A change in quantity demanded
            \choice A change in equilibrium only
        \end{choices}
    \end{mcquestion}

    \begin{mcquestion}[2]
        Utility maximization occurs when:
        \begin{choices}
            \choice Total utility is zero
            \CorrectChoice Marginal utility per dollar is equal across goods
            \choice Income is maximized
            \choice Prices are equal
        \end{choices}
    \end{mcquestion}

    \begin{mcquestion}[2]
        A deadweight loss occurs when:
        \begin{choices}
            \choice Markets are efficient
            \CorrectChoice Output is not at efficient equilibrium
            \choice Firms earn profit
            \choice Prices are stable
        \end{choices}
    \end{mcquestion}

    \begin{mcquestion}[2]
        In the long run, firms enter a market if:
        \begin{choices}
            \choice Profits are zero
            \CorrectChoice Economic profits are positive
            \choice Prices fall
            \choice Demand decreases
        \end{choices}
    \end{mcquestion}

    %%%%%%%%%%%%%%%%%%%%%%%%%%%%%%%%%%%%%%%%%%%%%%%%%%%%%%
    \section*{Part II: Short Answer (10 points each)}

    \begin{shortanswer}[10]{2in}
        Explain the difference between a change in demand and a change in quantity demanded. Include a simple example of each.
    \end{shortanswer}
    \begin{solution}
        A change in demand shifts the demand curve (income, tastes, etc.). A change in quantity demanded is movement along the curve due to price changes.
    \end{solution}

    \begin{shortanswer}[10]{2in}
        Draw (describe if necessary) what happens to equilibrium price and quantity when supply increases.
    \end{shortanswer}
    \begin{solution}
        Supply curve shifts right $\rightarrow$ price falls $\rightarrow$ quantity rises.
    \end{solution}

    \begin{shortanswer}[10]{2in}
        Explain why elasticity matters for determining the effect of a tax on consumers and producers.
    \end{shortanswer}
    \begin{solution}
        More inelastic side bears more tax burden because quantity adjusts less.
    \end{solution}

    \begin{shortanswer}[10]{2in}
        A firm experiences diminishing marginal returns. Explain what this means in production terms.
    \end{shortanswer}
    \begin{solution}
        Each additional worker adds less output than the previous worker.
    \end{solution}

    \begin{shortanswer}[10]{2in}
        Explain the relationship between marginal cost and average cost.
    \end{shortanswer}
    \begin{solution}
        MC below AC pulls AC down; MC above AC pushes AC up.
    \end{solution}

    \begin{shortanswer}[10]{2in}
        Why do perfectly competitive firms earn zero economic profit in the long run?
    \end{shortanswer}
    \begin{solution}
        Entry/exit adjusts supply until price equals minimum average cost.
    \end{solution}

    \begin{shortanswer}[10]{2in}
        Explain consumer surplus and how it is shown on a demand curve graph.
    \end{shortanswer}
    \begin{solution}
        Area between demand curve and price line.
    \end{solution}

    \begin{shortanswer}[10]{2in}
        Explain how a binding price ceiling creates shortages.
    \end{shortanswer}
    \begin{solution}
        Price below equilibrium increases demand and reduces supply.
    \end{solution}

    \begin{shortanswer}[10]{2in}
        A firm is deciding whether to shut down in the short run. What rule should it use?
    \end{shortanswer}
    \begin{solution}
        Shut down if price < AVC.
    \end{solution}

    \begin{shortanswer}[10]{2in}
        Explain one real-world example of opportunity cost in decision-making.
    \end{shortanswer}
    \begin{solution}
        Answers vary (time, money, foregone alternatives).
    \end{solution}

    %%%%%%%%%%%%%%%%%%%%%%%%%%%%%%%%%%%%%%%%%%%%%%%%%%%%%%
    \section*{Part III: Math Problems (10 points each)}

    \begin{problem}[10]{2in}
    Given demand $q = 100 - 2P$ and supply $q = 20 + 3P$, find equilibrium price and quantity.
    \end{problem}
    \begin{solution}
        $100 - 2P = 20 + 3P$ $\rightarrow$ $80 = 5P$ $\rightarrow$ $P = 16, q = 68$
    \end{solution}

    \begin{problem}[10]{2in}
    A product price rises from 10 to 12 and quantity demanded falls from 50 to 40. Compute price elasticity of demand.
    \end{problem}
    \begin{solution}
        % midpoint elasticity
        $E = \frac{(40-50)/45}{(12-10)/11} = \frac{-10/45}{2/11} = -2.44$
    \end{solution}

    \begin{problem}[10]{2in}
    A firm has cost function $TC = 50 + 10q + q^2$. Find marginal cost and average cost.
    \end{problem}
    \begin{solution}
        $MC = 10 + 2q,\quad AC = \frac{50}{q} + 10 + q$
    \end{solution}

    \begin{problem}[10]{2in}
    A tax of \$4 per unit is imposed on a market. Explain (no graph required) what happens to equilibrium price and quantity.
    \end{problem}
    \begin{solution}
        Supply shifts up; price rises, quantity falls.
    \end{solution}

    \begin{problem}[10]{3in}
    Given demand $P = 60 - Q$ and cost $TC = 20 + 10Q$, find profit-maximizing output and price.
    \end{problem}
    \begin{solution}
        $TR = (60-Q)Q = 60Q - Q^2$ $\rightarrow$ $MR = 60 - 2Q$
        $MC = 10$
        Set $MR = MC: 60 - 2Q = 10$ $\rightarrow$ $Q = 25$ $\rightarrow$ $P = 35$
    \end{solution}

    %%%%%%%%%%%%%%%%%%%%%%%%%%%%%%%%%%%%%%%%%%%%%%%%%%%%%%
    \section*{Bonus Question - 10 points}

    \begin{shortanswer}[10]{3in}
        What is one concept from microeconomics that you think helps explain real-world pricing or decision-making most clearly?
    \end{shortanswer}
    \begin{solution}
        Answers vary.
    \end{solution}

\end{questions}

\end{document}
